\scanpage{391}\mspage{\deleted{\unsure{129}}\unsure{130|129}}

\deleted{\unsure{ne se résoudra sans doute}}

Car je veux quand même que le type dans $\mathrm{Hot}$ soit le bon, qu'il corresponde bien à ``l'homotopie'' des espaces de chemins. Donc \deleted{que} on \unsure{voudra|voudrait} qu'un morphisme

\[
X\xrightarrow{f}X'
\]

dans $(\deleted{\unsure{\mathrm{Hot}}}\mathrm{Cat})$, qui est une $\mathrm{Hot}$-équivalence, définisse des

\[
\underline{\mathrm{Ch}}_\infty(X\unsure{,|;}a,b)\unsure{-|\longrightarrow}\underline{\mathrm{Ch}}_\infty(X',f(a),f(\unsure{a|b}))
\]

qui soient eux aussi des $\mathrm{Hot}$-équivalences.

D'autre part, c'est évidemment compatible avec la \unsure{brillante} composition des chemins, que j'ai réussi à définir \unsure{entre} ces objets de $\mathrm{Hot}$. Donc \deleted{pour} (sous réserve d'avoir vérifié ce que je viens de dire) la vérification de l'associativité se ramènerait au cas où $X$ est de Thomason-Kan. Mais dans ce cas, il y a des chances que, tout comme dans le cas $1/2$ simplicial, les $\underline{\mathrm{Ch}}_\infty(X;a,b)$ peuvent se définir \deleted{\unsure{en ter}} (comme objets de $\mathrm{Hot}$) comme étant simplement les chemins de type $\Delta_+$ de $a$ \unsure{à} $b$, qu'on compose alors de \unsure{façon} triviale. Mais non, ces catégories sont discrètes !

Je renonce pour l'instant à poursuivre \deleted{\unsure{plus}} la question d'associativité. Demander à Thomason pour les propriétés du foncteur $\mathrm{Hot}\longrightarrow\widehat{\mathrm{Hot}_0}$.

% archive sheet number: 390
